\documentclass[leqno]{article}
% Shared mark scheme preamble for resources.danieljsmith.org: the pack
% preamble, unchanged, then the scheme's own layer. A scheme's .tex is
% rendered from its YAML (scripts/README.md, Mark schemes), never edited.
\input{\CurrentFilePath/pack_preamble.tex}

% The page: equal side margins, so the table runs the full width of the
% header rule. The footers' right edge follows the text's.
\geometry{left=0.8in, right=0.8in}
\makeatletter
\renewcommand{\djs@footR}[1]{%
  \put(\dimexpr1in+\hoffset+\oddsidemargin+\textwidth\relax,
    -\dimexpr\paperheight-0.32in\relax){\makebox[0pt][r]{#1}}}
\makeatother

\usepackage[export]{adjustbox}
\usepackage{tabularray}
\definecolor{msGrid}{RGB}{176, 172, 164}
\definecolor{msTint}{RGB}{245, 243, 239}

% The header on every page, the front page included, and the booklet label
% under the contents-page title. \djsQuestionPageStyle stays empty.
\newcommand{\djsSchemeHeader}[1]{%
  \djsPackHeader[markscheme]{\textit{Mark scheme} - #1}%
  \renewcommand{\djsPackTitle}{#1}%
  \renewcommand{\djsBookletLabel}{Mark scheme}%
}

% The front page: the contents heading, spaced as tocloft spaces the packs'
% "Questions", then the contents table. A question's name and page number
% both link to its first page.
\newcommand{\msContentsHeading}{%
  \vskip\cftbeforetoctitleskip
  {\noindent\cfttoctitlefont\contentsname\par}%
  \vskip\cftaftertoctitleskip}
\NewTblrEnviron{mscontents}
\SetTblrInner[mscontents]{
  colspec = {Q[wd=26mm, l] X[l] Q[wd=16mm, c] Q[wd=14mm, r]},
  width = \linewidth,
  rowsep = 4pt, colsep = 6pt,
  hlines = {0.4pt, msGrid}, vlines = {0.4pt, msGrid},
  hline{2} = {2.2pt, scarlet},
  row{1} = {font=\small\bfseries},
  row{Z} = {font=\bfseries},
  column{2} = {font=\small},
}
\newcommand{\msContentsQuestion}[1]{\hyperref[q:#1]{Question #1}}
\newcommand{\msContentsPage}[1]{\hyperref[q:#1]{\pageref*{q:#1}}}
\newcommand{\msSep}{\hspace{0.9em}}

% The conventions pages' layout: headings in the documents' own bold, tables
% ruled and padded as the scheme's, roomy lists.
\newcommand{\msHead}[1]{%
  \par\vspace{18pt}%
  {\noindent\bfseries #1\par}%
  \vspace{6pt}}
\NewTblrEnviron{mstblr}
\SetTblrInner[mstblr]{rowsep = 6pt, colsep = 6pt,
  hlines = {0.4pt, msGrid}, vlines = {0.4pt, msGrid}}
\newenvironment{msList}
  {\begin{itemize}[leftmargin=1.2em, itemsep=5pt, topsep=3pt, parsep=0pt]}
  {\end{itemize}}

% A question's first page: the link target and the PDF bookmark, as
% \questionitem gives the packs.
\newcommand{\msQnum}{}
\newcommand{\msQuestion}[1]{%
  \clearpage
  \renewcommand{\msQnum}{#1}%
  \phantomsection\label{q:#1}%
  \addcontentsline{toc}{subsection}{\protect\djsTocEntry{#1}}}

% The scheme's table: Question 16mm, Scheme the rest, Marks 16mm, Notes
% 44mm; every cell ruled in msGrid; the scarlet rule under the head row. A
% long table has no caption or "continued" foot; a continuation page opens
% with a muted line above the repeated head row. The rules across a part's
% rows, the alternative's tint and the total row are the rendered table's.
\NewTblrTheme{ms}{%
  \DefTblrTemplate{firsthead}{default}{}%
  \DefTblrTemplate{middlehead,lasthead}{default}{%
    {\small\itshape\color{djsMuted}Question \msQnum\ continued\par}%
    \vspace{6pt}}%
  \DefTblrTemplate{firstfoot,middlefoot,lastfoot}{default}{}%
}
\NewTblrEnviron{msscheme}
\SetTblrOuter[msscheme]{long, theme=ms, entry=none, label=none}
\SetTblrInner[msscheme]{
  colspec = {Q[wd=16mm, l] X[l] Q[wd=16mm, c] Q[wd=44mm, l]},
  width = \linewidth,
  rowsep = 6pt, colsep = 6pt,
  rowhead = 1,
  vlines = {0.4pt, msGrid},
  hline{1} = {0.4pt, msGrid}, hline{2} = {2.2pt, scarlet},
  hline{Z} = {0.4pt, msGrid},
  column{1} = {font=\bfseries}, column{3} = {font=\bfseries},
  column{4} = {font=\small},
  row{1} = {rowsep=4pt, font=\small\bfseries},
}

% Two marks on a row: each note starts with its code.
\newcommand{\msNote}[2]{\textbf{#1}\enspace #2}
% A full-row figure: the worked solution's own picture, scaled uniformly and
% only down, to the table less the Question column. The picture is drawn
% with the line width it has in the solutions document, at its list depth
% there (the optional argument, 1 by default: a part's working), so one
% sized by \linewidth keeps its size; tabularray sets the cell's \linewidth
% to the cell's measured width, which would shrink it at every pass. The
% widths are measured once, in a solution box at the packs' text width
% (pack_preamble.tex's margins, 0.8in and 1in).
\newlength{\msSpanW}
\makeatletter
\newcommand{\ms@keepLW}[1]{%
  \expandafter\xdef\csname ms@solnLW@#1\endcsname{\the\linewidth}}
\AtBeginDocument{%
  \setlength{\msSpanW}{\dimexpr\textwidth-16mm-28pt\relax}%
  \setbox\z@\vbox{%
    \begin{minipage}{\dimexpr\paperwidth-1.8in\relax}
      \begin{djsSolution}
        \ms@keepLW{0}%
        \begin{enumerate}\item \ms@keepLW{1}%
        \begin{enumerate}\item \ms@keepLW{2}%
        \begin{enumerate}\item \ms@keepLW{3}%
        \end{enumerate}\end{enumerate}\end{enumerate}
      \end{djsSolution}
    \end{minipage}}}
\newcommand{\msFigure}[2][1]{%
  \adjustbox{max width=\msSpanW, max totalheight=0.6\textheight, valign=t}{%
    \hsize=\csname ms@solnLW@#1\endcsname
    \linewidth=\hsize
    #2}}
\makeatother
% A figure mark's feature, in the Scheme column.
\newcommand{\msFeature}[1]{\textit{#1}}
% An alternative's title row.
\newcommand{\msAlternative}[1]{\textit{Alternative:} #1}
% Drafts only: a queried mark's code, the query's why under its note, and a
% flag in its part's place.
\newcommand{\msQueried}[2]{#1~\(\to\) #2?}
\newcommand{\msQuery}[1]{{\itshape\color{djsMuted}Query: #1}}
\newcommand{\msFlag}[2]{\textbf{#1:}\enspace #2}


\djsSchemeHeader{Polar Coordinates Basics}
\djsSitePath{Core Pure}{Polar Coordinates}

\djsLastUpdated{25 October 2026}
\begin{document}
\djsSetQuestionCount{10}
\djsAddMarks{1}{8}
\djsAddMarks{2}{8}
\djsAddMarks{3}{7}
\djsAddMarks{4}{8}
\djsAddMarks{5}{11}
\djsAddMarks{6}{7}
\djsAddMarks{7}{12}
\djsAddMarks{8}{13}
\djsAddMarks{9}{8}
\djsAddMarks{10}{15}
\vspace*{6pt}
\djsContentsTitle
\vspace{30pt}
\msContentsHeading
\begin{mscontents}{}
Question & Marks by part & Marks & Page \\
\msContentsQuestion{1} & (a)\,3\msSep (b)\,2\msSep (c)\,3 & 8 & \msContentsPage{1} \\
\msContentsQuestion{2} & (a)\,1\msSep (b)\,3\msSep (c)\,4 & 8 & \msContentsPage{2} \\
\msContentsQuestion{3} & (a)\,1\msSep (b)\,6 & 7 & \msContentsPage{3} \\
\msContentsQuestion{4} & (a)\,4\msSep (b)\,4 & 8 & \msContentsPage{4} \\
\msContentsQuestion{5} & (a)\,4\msSep (b)\,4\msSep (c)\,3 & 11 & \msContentsPage{5} \\
\msContentsQuestion{6} & (i)\,4\msSep (ii)\,3 & 7 & \msContentsPage{6} \\
\msContentsQuestion{7} & (a)\,3\msSep (b)\,3\msSep (c)\,2\msSep (d)\,4 & 12 & \msContentsPage{7} \\
\msContentsQuestion{8} & (i)\,2\msSep (ii)(a)\,4\msSep (ii)(b)\,3\msSep (iii)\,4 & 13 & \msContentsPage{8} \\
\msContentsQuestion{9} & (a)\,2\msSep (b)\,6 & 8 & \msContentsPage{9} \\
\msContentsQuestion{10} & (a)\,3\msSep (b)(i)\,3\msSep (b)(ii)\,5\msSep (b)(iii)\,4 & 15 & \msContentsPage{10} \\
\SetCell[c=2]{r} Total & & 97 & \\
\end{mscontents}
\newpage
% The conventions pages, shared by every pack's mark scheme: one source of
% truth, \input after the front page, with no strand or pack logic. The keys
% on the first side, the rules on the second. The scheme layer defines
% \msHead, msList and mstblr.
\vspace*{6pt}%
{\noindent\cfttoctitlefont Reading the mark scheme\par}%
\vspace{10pt}%
These rules apply throughout, unless a note says otherwise.

\msHead{Marks}
Every mark is worth one, so \textbf{M1 A1} on a row is two marks. Only marks
coded ft follow through.\par
\vspace{10pt}%
\begin{mstblr}{
  colspec = {Q[wd=15mm, l] X[l]},
  width = \linewidth,
  hline{2} = {2.2pt, scarlet},
  column{1} = {font=\bfseries},
  row{1} = {font=\small\bfseries, rowsep=4pt},
}
Code & Meaning \\
M1 & Method: a correct method, applied to this question. A slip doesn't cost it if the method is clear. A formula must be used, not just quoted. \\
dM1 & Dependent method: as M1, but only if the method mark before it was earned, unless its note names another. \\
A1 & Accuracy: a correct answer or step. It needs the method mark it follows. \\
A1ft & As A1, following through from an earlier answer that the note names. \\
A1~AG & The last mark of a part whose answer is printed in the question: complete working, with no errors. \\
B1 & A correct value, fact or statement. It needs no method mark. \\
B1ft & As B1, following through from an earlier answer that the note names. \\
E1 & A written explanation, interpretation or conclusion. Concluding a proof by induction needs every earlier mark. \\
\end{mstblr}

\msHead{Abbreviations}
ft and AG qualify a code in the Marks column. The others appear in the
notes.\par
\vspace{10pt}%
\begin{mstblr}{
  colspec = {Q[wd=17mm, l] X[l] Q[wd=15mm, l] X[l]},
  width = \linewidth,
  column{1} = {font=\bfseries}, column{3} = {font=\bfseries},
}
ft & follow through & AG & answer given in the question \\
cao & correct answer only, with no follow-through & cso & correct solution only: no errors in the part \\
oe & or an equivalent form & awrt & any answer that rounds to this value \\
soi & seen, or implied by later working & isw & ignore working after a correct answer \\
SC & special case: credit outside the main method & their & the student's own earlier value \\
condone & a slip that doesn't cost the mark & [\;\,] & working the mark doesn't need \\
s.f. & significant figures & d.p. & decimal places \\
\end{mstblr}

\newpage
\msHead{Reading the scheme}
\begin{msList}
\item Each row is a line of working, with the marks it earns beside it.
  Where a row carries two marks, each note starts with its code.
\item \textbf{[\textit{n}]} ends each part with its total.
\item A shaded \textit{Alternative} is a second method for the same total.
  Mark an answer by the method it follows.
\item Any other valid method earns the equivalent marks, unless the question
  names a method or says \textit{hence}.
\item A sketch earns one B1 for each feature named below its figure.
\end{msList}

\msHead{Accuracy and answers}
\begin{msList}
\item Unless a question says otherwise, give answers to 3 significant
  figures, and angles in degrees to 1 decimal place. Angles in radians take
  3 significant figures.
\item A final answer with more figures is accepted if it rounds to the
  scheme's value; one with fewer is not. An answer spoiled by rounding too
  early loses its accuracy mark.
\item Where a question asks for an exact answer, a decimal loses the final
  accuracy mark. Otherwise either is accepted.
\item Wrong or missing units cost nothing unless the question asks for them.
\item A correct answer with no working earns full marks, unless the question
  needs the working: show that, prove, hence, a method it names, or an exact
  answer. A wrong answer with no working earns nothing. Calculator answers
  are treated the same way.
\item A correct answer, once reached, isn't spoiled by later working (isw).
\end{msList}

\msHead{Value of $\boldsymbol{g}$}
\begin{msList}
\item Take $g = 9.8\,\mathrm{m\,s^{-2}}$ unless a question states a value,
  which is then the only one accepted. Answers that use $g = 9.8$ may be
  given to 2 or 3 significant figures.
\item Where a question doesn't state $g$, an answer worked with $g = 10$
  throughout earns the same marks, keeping any value printed in the
  question. The worked solutions list the final answers with $g = 10$.
\end{msList}


\msQuestion{1}
\begin{msscheme}{
  hline{3-5,7-8,10-13} = {2-Z}{0.4pt, msGrid},
  hline{6,9,14} = {0.4pt, msGrid},
  row{14} = {font=\bfseries},
}
Question & Scheme & Marks & Notes \\
1(a) & {\(A = (1+2\pi, \pi)\)} & B1 & {Accept \(r\) awrt 7.28; \(\theta = \pi\) or \(180^\circ\).} \\
& {\(x = \their{(1+2\pi)}\cos\pi,\ y = \their{(1+2\pi)}\sin\pi\)} & M1 & {Uses \(x = r\cos\theta\) and \(y = r\sin\theta\) with their \(r\) and \(\theta = \pi\).} \\
& {\((-1-2\pi, 0)\)} & A1 & {oe, or awrt \((-7.28, 0)\).} \\
\nopagebreak & & [3] & \\
1(b) & {\(1+2\theta = 8-\theta \Rightarrow \theta = \ldots\)} & M1 & {Equates the two expressions for \(r\) and solves for \(\theta\). Condone sign slips.} \\
& {\(\left(\frac{17}{3}, \frac73\right)\)} & A1 & {oe, or awrt \((5.67, 2.33)\). Both coordinates needed.} \\
\nopagebreak & & [2] & \\
1(c) & \SetCell[c=3]{c} \msFigure{\begin{tikzpicture}[scale=0.75]
\def\rOuter{8}
\def\rStart{1}
\def\growth{2}
\pgfmathsetmacro{\thOne}{pi}
\pgfmathsetmacro{\thTwo}{2*pi}
\pgfmathsetmacro{\rInner}{\rOuter-\thTwo}
\pgfmathsetmacro{\rA}{\rStart+\growth*\thOne}
\pgfmathsetmacro{\axisEnd}{1.15*\rOuter}
\pgfmathsetmacro{\labelOne}{\rStart+\growth*pi/2}
\pgfmathsetmacro{\labelTwo}{\rOuter-pi/2}
\coordinate (O) at (0,0);
\coordinate (A) at (180:\rA);
\coordinate (S) at (0:\rOuter);
\coordinate (T) at (0:\rInner);
\coordinate (L) at (\axisEnd,0);
\coordinate (LabelOne) at (90:\labelOne);
\coordinate (LabelTwo) at (90:\labelTwo);
\draw[thick, black, samples=150, domain=0:\thOne, variable=\t]
  plot ({(\rStart+\growth*\t)*cos(deg(\t))},{(\rStart+\growth*\t)*sin(deg(\t))});
\draw[thick, black, samples=200, domain=0:\thTwo, variable=\t]
  plot ({(\rOuter-\t)*cos(deg(\t))},{(\rOuter-\t)*sin(deg(\t))});
\draw[-{Stealth[length=3mm]}, thick, black] (O) -- (L);
\node[below left=2pt] at (O) {$O$};
\node[left=2pt] at (A) {$A$};
\node[below right] at (L) {Initial line};
\node[below left] at (LabelOne) {$C_1$};
\node[above right] at (LabelTwo) {$C_2$};
\foreach \P in {S,T} {\fill[black] (\P) circle (1.5pt);}
\node[above right=2pt] at (S) {$(8,0)$};
\node[below right=2pt] at (T) {$(8-2\pi,2\pi)$};
\end{tikzpicture}} & & \\
& \msFeature{Shape} & B1 & {Spiral making one anticlockwise turn with \(r\) decreasing throughout, starting and ending on the initial line, not reaching \(O\).} \\
& \msFeature{Labelled point} & B1 & {Start on the initial line labelled 8 or \((8, 0)\).} \\
& \msFeature{Labelled point} & B1 & {End on the initial line between \(O\) and the start, labelled \(8-2\pi\), awrt 1.72, or \((8-2\pi, 2\pi)\).} \\
\nopagebreak & & [3] & \\
\nopagebreak \SetCell[c=2]{r} Total for Question 1 & & 8 & \\
\end{msscheme}

\msQuestion{2}
\begin{msscheme}{
  hline{3,5-7,9-12} = {2-Z}{0.4pt, msGrid},
  hline{4,8,13} = {0.4pt, msGrid},
  row{13} = {font=\bfseries},
}
Question & Scheme & Marks & Notes \\
2(a) & {\(r = \dfrac{6}{2-\sin\frac{\pi}{6}} = \dfrac{6}{2-\frac12} = 4\)} & B1 & {Needs \(\sin\frac{\pi}{6} = \frac12\) substituted, or the denominator \(\frac32\) seen, leading to \(r = 4\).} \\
\nopagebreak & & [1] & \\
2(b) & {\(r\) greatest when \(2-\sin\theta\) is least: \(\sin\theta = 1\)} & B1 & {Identifies \(\sin\theta = 1\), or \(\theta = \frac{\pi}{2}\), as giving greatest \(r\), by the least denominator or otherwise.} \\
& {\(\theta = \frac{\pi}{2} \Rightarrow r = \dfrac{6}{2-1} = \ldots\)} & M1 & {Substitutes their \(\theta\), from maximising \(r\), into the equation of \(C\) to find \(r\).} \\
& {\(\left(6, \frac{\pi}{2}\right)\)} & A1 & {Accept \(r = 6\), \(\theta = \frac{\pi}{2}\) stated separately.} \\
\nopagebreak & & [3] & \\
2(c) & {\(r\sin\theta = r\cos\theta \Rightarrow \tan\theta = 1\)} & M1 & {Substitutes \(x = r\cos\theta\), \(y = r\sin\theta\) into \(y = x\), reaching \(\tan\theta = 1\) or an angle directly.} \\
& {\(\theta = \frac{\pi}{4},\ -\frac{3\pi}{4}\)} & A1 & {Both, with no others in \(-\pi<\theta\leq\pi\). Condone \(\frac{5\pi}{4}\) for the second.} \\
& {\(r = \dfrac{6}{2-\frac{\sqrt2}{2}},\ x = r\cos\frac{\pi}{4} = \ldots\)\\[4pt] \(r = \dfrac{6}{2+\frac{\sqrt2}{2}},\ x = r\cos\left(-\frac{3\pi}{4}\right) = \ldots\)} & M1 & {Finds \(r\) at one of their angles and uses \(x = r\cos\theta\) or \(y = r\sin\theta\). Need not be simplified.} \\
& {\(\left(\frac{6+12\sqrt2}{7}, \frac{6+12\sqrt2}{7}\right)\), \(\left(\frac{6-12\sqrt2}{7}, \frac{6-12\sqrt2}{7}\right)\)} & A1 & {Both points, each with both coordinates, in any exact form; need not be rationalised.} \\
\nopagebreak & & [4] & \\
\nopagebreak \SetCell[c=2]{r} Total for Question 2 & & 8 & \\
\end{msscheme}

\msQuestion{3}
\begin{msscheme}{
  hline{3-4,6-12} = {2-Z}{0.4pt, msGrid},
  hline{5,13} = {0.4pt, msGrid},
  row{13} = {font=\bfseries},
}
Question & Scheme & Marks & Notes \\
3(a) & \SetCell[c=3]{c} \msFigure{\begin{tikzpicture}[scale=0.85]
\def\base{4}
\def\amp{3}
\def\axislen{2.5}
\def\leftspace{7.5}
\def\rightspace{3.5}
\def\verticalspace{5.3}
\coordinate (O) at (0,0);
\coordinate (InitialEnd) at (\axislen,0);
\path[use as bounding box] (-\leftspace,-\verticalspace) rectangle (\rightspace,\verticalspace);
\draw[thick, black, domain=-180:180, samples=300, variable=\t]
  plot ({(\base-\amp*cos(\t))*cos(\t)},{(\base-\amp*cos(\t))*sin(\t)});
\draw[-{Stealth[length=3mm]}, thick, black] (O) -- (InitialEnd)
  node[right] {Initial line};
\node[below left=2pt] at (O) {$O$};
\end{tikzpicture}} & & \\
& \msFeature{Symmetry} & B1 & {Given portion reflected in the initial line, giving a closed curve symmetric about it and still passing through \((1,0)\).} \\
\nopagebreak & & [1] & \\
3(b) & \SetCell[c=3]{c} \msFigure{\begin{tikzpicture}[scale=0.85]
\def\base{4}
\def\amp{3}
\def\axislen{2.5}
\def\leftspace{7.5}
\def\rightspace{3.5}
\def\verticalspace{5.3}
\def\lineX{1}
\pgfmathsetmacro{\interR}{\base-\amp/3}
\pgfmathsetmacro{\interAng}{acos(1/3)}
\coordinate (O) at (0,0);
\coordinate (InitialEnd) at (\axislen,0);
\coordinate (LineBottom) at (\lineX,-\verticalspace);
\coordinate (LineTop) at (\lineX,\verticalspace);
\coordinate (IntersectionZero) at (\lineX,0);
\coordinate (IntersectionUpper) at (\interAng:\interR);
\coordinate (IntersectionLower) at (-\interAng:\interR);
\path[use as bounding box] (-\leftspace,-\verticalspace) rectangle (\rightspace,\verticalspace);
\draw[thick, black, domain=-180:180, samples=300, variable=\t]
  plot ({(\base-\amp*cos(\t))*cos(\t)},{(\base-\amp*cos(\t))*sin(\t)});
\draw[thick, black] (LineBottom) -- (LineTop) node[above right] {$\ell$};
\draw[-{Stealth[length=3mm]}, thick, black] (O) -- (InitialEnd)
  node[right] {Initial line};
\node[below left=2pt] at (O) {$O$};
\foreach \P in {IntersectionZero,IntersectionUpper,IntersectionLower}
  {\fill[black] (\P) circle (1.5pt);}
\node[above right=3pt] at (IntersectionZero) {$(1,0)$};
\node[above right=3pt] at (IntersectionUpper) {$(3,1.23)$};
\node[below right=3pt] at (IntersectionLower) {$(3,-1.23)$};
\end{tikzpicture}} & & \\
& \msFeature{Intercepts} & B1 & {Vertical line through \((1,0)\) on the initial line, meeting \(C\) there and crossing it once above and once below the initial line.} \\
& {\((4-3\cos\theta)\cos\theta = 1\)} & M1 & {Substitutes for \(r\) in \(r\cos\theta=1\) to obtain an equation in \(\cos\theta\) alone.} \\
& {\(3\cos^2\theta - 4\cos\theta + 1 = 0\)} & A1 & {oe, any three-term form.} \\
& {\((3\cos\theta-1)(\cos\theta-1)=0 \Rightarrow \cos\theta = \ldots\)\\[4pt] \(\theta = \ldots,\ r = 4-3\cos\theta = \ldots\)} & dM1 & {Solves their three-term quadratic in \(\cos\theta\) by any means, then finds at least one \(\theta\) and its \(r\) from either equation.} \\
& {\((1, 0)\)} & A1 & \\
& {\((3, 1.23)\), \((3, -1.23)\)} & A1 & {Both, with \(r = 3\) and angles awrt \(\pm1.23\). A0 for an angle outside \(-\pi<\theta\leq\pi\).} \\
\nopagebreak & & [6] & \\
\nopagebreak \SetCell[c=2]{r} Total for Question 3 & & 7 & \\
\end{msscheme}

\msQuestion{4}
\begin{msscheme}{
  hline{3-5,7-11} = {2-Z}{0.4pt, msGrid},
  hline{6,12} = {0.4pt, msGrid},
  row{12} = {font=\bfseries},
}
Question & Scheme & Marks & Notes \\
4(a) & {\(x^2+y^2+4x-6y=12\)} & M1 A1 & {\msNote{M1}{Replaces \(r^2\) by \(x^2+y^2\) and at least one of \(r\cos\theta\), \(r\sin\theta\) by \(x\), \(y\).}\\ \msNote{A1}{Correct equation, oe.}} \\
& {\((x+2)^2-4+(y-3)^2-9=12\)} & M1 & {Completes the square in both \(x\) and \(y\) for their equation of the form \(x^2+y^2+ax+by=c\). Condone sign slips.} \\
& {\((x+2)^2+(y-3)^2=25\)\\[4pt] Centre \((-2,3)\), radius \(5\)} & A1 & {Both centre and radius, from a correct equation in circle form.} \\
\nopagebreak & & [4] & \\
4(b) & \SetCell[c=3]{c} \msFigure{\begin{tikzpicture}
\begin{axis}[
width=0.8\linewidth,
height=8cm,
axis lines=middle,
axis equal image,
xmin=-8, xmax=4,
ymin=-3, ymax=9,
xtick=\empty,
ytick=\empty,
xlabel={$x$},
ylabel={$y$},
clip=false
]
\addplot[domain=0:360,samples=181,smooth] ({-2+5*cos(x)},{3+5*sin(x)});
\fill (axis cs:-2,3) circle (2pt);
\node[anchor=east] at (axis cs:-2.2,3) {\((-2,3)\)};
\fill (axis cs:0,0) circle (2pt);
\node[anchor=north east] at (axis cs:-0.15,-0.15) {\(O(0,0)\)};
\node at (axis cs:2.4,6.6) {\(C_1\)};
\end{axis}
\end{tikzpicture}} & & \\
& \msFeature{Centre and radius} & B1 & {Circle with its centre marked in the second quadrant, labelled \((-2,3)\).} \\
& \msFeature{Labelled point} & B1 & {Pole marked at the origin, inside the circle.} \\
& {\(OA=\sqrt{(-2)^2+3^2}\ [=\sqrt{13}]\)\\[4pt] \(5\pm\their{\sqrt{13}}\)} & M1 & {Finds the distance from the pole to their centre and adds it to and subtracts it from their radius. No justification that the pole is inside is needed.} \\
& {Greatest \(5+\sqrt{13}\), least \(5-\sqrt{13}\)} & A1 & {Both, correctly identified.} \\
\nopagebreak & & [4] & \\
\nopagebreak \SetCell[c=2]{r} Total for Question 4 & & 8 & \\
\end{msscheme}

\msQuestion{5}
\begin{msscheme}{
  hline{3-6,8-11,13-16} = {2-Z}{0.4pt, msGrid},
  hline{7,12,17} = {0.4pt, msGrid},
  row{17} = {font=\bfseries},
}
Question & Scheme & Marks & Notes \\
5(a) & {\(x^2+\left(y-\frac{\sqrt3}{2}\right)^2=\ldots\)} & M1 & {Completes the square in \(y\), or compares with the general circle equation. Condone a slip in the constant.} \\
& {Centre \(\left(0,\frac{\sqrt3}{2}\right)\), radius \(\frac{\sqrt3}{2}\)} & A1 & {Both, exact or awrt 0.866.} \\
& {\(r^2\cos^2\theta+r^2\sin^2\theta=\sqrt3\,r\sin\theta\)} & M1 & {Substitutes \(x=r\cos\theta\), \(y=r\sin\theta\), or uses \(x^2+y^2=r^2\) with \(y=r\sin\theta\).} \\
& {\(r^2=\sqrt3\,r\sin\theta\), \(r\neq0\)\\[4pt] \(\Rightarrow r=\sqrt3\sin\theta\)} & A1~AG & {Correct substituted equation leading without error to the printed result. Needs \(r\neq0\), or non-pole, stated when dividing by \(r\).} \\
\nopagebreak & & [4] & \\
5(b) & {\(\sqrt3\sin\theta=1+\cos\theta\)\\[4pt] \(\Rightarrow \sqrt3\sin\theta-\cos\theta=R\sin(\theta-\alpha)=1\)} & M1 & {Equates the radii and reduces to one trigonometric function, e.g. any harmonic form with \(R\) and \(\alpha\) attempted, or half-angle identities. Condone sign slips.} \\
& {\(R=2,\ \alpha=\frac{\pi}{6}\)} & A1 & {Or a correct equivalent single-function equation by another route.} \\
& {\(\sin\left(\theta-\frac{\pi}{6}\right)=\frac12 \Rightarrow \theta=\ldots\)} & dM1 & {Solves their equation for \(\theta\) in \(0\le\theta\le\pi\), finding at least one value.} \\
& {\(\theta=\frac{\pi}{3}\) (\(\theta=\pi\) is the pole)\\[4pt] \(P\left(\frac32,\frac{\pi}{3}\right)\)} & A1 & {cso. Exact. Needs \(\theta=\pi\) rejected, or \(\frac{\pi}{3}\) shown to be the only solution in \(0<\theta<\pi\), with one intersection stated.} \\
\nopagebreak & & [4] & \\
5(c) & \SetCell[c=3]{c} \msFigure{\begin{tikzpicture}[x=3.2cm,y=3.2cm]
\draw[->] (-1.2,0) -- (2.5,0) node[right] {\(x\)};
\draw[->] (0,-0.3) -- (0,2.1) node[above] {\(y\)};
\draw[domain=0:180,samples=181,smooth,variable=\t]
  plot ({sqrt(3)*sin(\t)*cos(\t)},{sqrt(3)*sin(\t)*sin(\t)});
\draw[domain=0:180,samples=181,smooth,variable=\t]
  plot ({(1+cos(\t))*cos(\t)},{(1+cos(\t))*sin(\t)});
\node at (-0.85,1.55) {\(C_1\)};
\node at (1.95,0.8) {\(C_2\)};
\fill (0,0) circle (1.5pt);
\node[below left] at (0,0) {\((0,0)\)};
\fill (2,0) circle (1.5pt);
\node[below] at (2,0) {\((2,0)\)};
\fill (0.75,{3*sqrt(3)/4}) circle (1.5pt);
\node[above right] at (0.75,{3*sqrt(3)/4}) {\(P\left(\frac32,\frac{\pi}{3}\right)\)};
\end{tikzpicture}} & & \\
& \msFeature{Shape} & B1 & {Complete circle above the initial line, touching it at the pole.} \\
& \msFeature{Shape} & B1 & {Upper half of a cardioid from \((2,0)\), passing into \(x<0\) before reaching the pole; both curves labelled.} \\
& \msFeature{Intercepts} & B1 & {\(P\) marked at the intersection, and \((2,0)\) and \((0,0)\) indicated.} \\
\nopagebreak & & [3] & \\
\nopagebreak \SetCell[c=2]{r} Total for Question 5 & & 11 & \\
\end{msscheme}

\msQuestion{6}
\begin{msscheme}{
  hline{3-6,8-10} = {2-Z}{0.4pt, msGrid},
  hline{7,11} = {0.4pt, msGrid},
  row{11} = {font=\bfseries},
}
Question & Scheme & Marks & Notes \\
6(i) & {\(r^4 = 5r^2\cos^2\theta + 4r^2\cos\theta\sin\theta\)\\[4pt] \(\phantom{r^4 =} + 2r^2\sin^2\theta\)} & M1 & {Substitutes \(x=r\cos\theta\), \(y=r\sin\theta\) and \(x^2+y^2=r^2\) throughout. Condone one slip.} \\
& {\(r^2 = 5\cos^2\theta + 4\cos\theta\sin\theta + 2\sin^2\theta\)} & A1 & \\
& {\(r^2 = \tfrac52(1+\cos2\theta) + 2\sin2\theta\)\\[4pt] \(\phantom{r^2 =} + (1-\cos2\theta)\)} & M1 & {Uses double-angle identities for all of \(\cos^2\theta\), \(\sin^2\theta\) and \(\sin\theta\cos\theta\). Condone sign and coefficient slips.} \\
& {\(r^2 = \dfrac{7+3\cos2\theta+4\sin2\theta}{2}\)} & A1~AG & {Correct equation in \(r\) and \(\theta\) leading without error to the printed result; combined steps accepted.} \\
\nopagebreak & & [4] & \\
6(ii) & {\(3\cos2\theta+4\sin2\theta = R\cos(2\theta-\alpha)\)\\[4pt] \(\dfrac{7-R}{2} \leq r^2 \leq \dfrac{7+R}{2}\)} & M1 & {Expresses \(3\cos2\theta+4\sin2\theta\) as any single harmonic form and uses its bounds \(\pm R\) in \(r^2\). Condone a wrong \(R\).} \\
& {\(R = 5 \Rightarrow 1 \leq r^2 \leq 6\)} & A1 & \\
& {Least distance \(1\), greatest distance \(\sqrt6\)} & A1 & {Both distances, exact. Not for the values of \(r^2\) alone.} \\
\nopagebreak & & [3] & \\
\nopagebreak \SetCell[c=2]{r} Total for Question 6 & & 7 & \\
\end{msscheme}

\msQuestion{7}
\begin{msscheme}{
  hline{3-5,7-10,12-13,15-17} = {2-Z}{0.4pt, msGrid},
  hline{6,11,14,18} = {0.4pt, msGrid},
  row{18} = {font=\bfseries},
}
Question & Scheme & Marks & Notes \\
7(a) & {\(r\cos\left(\theta+\tfrac{\pi}{6}\right)=6\)\\[4pt] \(r\left(\cos\theta\cos\tfrac{\pi}{6}-\sin\theta\sin\tfrac{\pi}{6}\right)=6\)} & M1 & {Multiplies through by \(\cos\left(\theta+\frac{\pi}{6}\right)\) and expands the compound angle. Condone sign slips in the expansion.} \\
& {\(\tfrac{\sqrt3}{2}x-\tfrac12y=6\)} & dM1 & {Replaces \(r\cos\theta\) by \(x\) and \(r\sin\theta\) by \(y\).} \\
& {\(y=\sqrt3x-12\)} & A1 & {cso. Must be in the form \(y=mx+c\); another linear form such as \(\sqrt3x-y=12\) scores A0.} \\
\nopagebreak & & [3] & \\
7(b) & \SetCell[c=3]{c} \msFigure{\begin{tikzpicture}[scale=0.5]
\def\sketchangle{40}
\def\initiallen{15}
\def\raylen{12}
\def\lowerheight{-12}
\def\upperheight{6}
\pgfmathsetmacro{\initialintercept}{4*sqrt(3)}
\pgfmathsetmacro{\lowerx}{(\lowerheight+12)/sqrt(3)}
\pgfmathsetmacro{\upperx}{(\upperheight+12)/sqrt(3)}
\coordinate (O) at (0,0);
\coordinate (E) at (\initiallen,0);
\coordinate (F) at (-\sketchangle:\raylen);
\coordinate (A) at (\initialintercept,0);
\coordinate (L) at (\lowerx,\lowerheight);
\coordinate (U) at (\upperx,\upperheight);
\draw[thick, black, -{Stealth[length=3mm]}] (O) -- (E)
  node[above left] {Initial line};
\draw[thick, black] (L) -- (U) node[pos=0.93, above left] {$G_1$};
\draw[thick, black, -{Stealth[length=3mm]}] (O) -- (F)
  node[below right] {$G_2$};
\pic["$\alpha$", draw, angle radius=8mm, angle eccentricity=1.4]
  {angle=F--O--E};
\fill[black] (A) circle (2pt);
\node[above left=2pt] at (A) {$\left(4\sqrt3,0\right)$};
\node[below left=2pt] at (O) {$O$};
\end{tikzpicture}} & & \\
& \msFeature{Shape} & B1 & {Straight line with positive gradient crossing the initial line to the right of \(O\), extending above and below it.} \\
& \msFeature{Half-line} & B1 & {Ray from \(O\) below the initial line meeting \(G_1\), with \(\alpha\) marked between the ray and the initial line.} \\
& \msFeature{Labelled point} & B1 & {\(\left(4\sqrt3,0\right)\), \(\left(6\sec\frac{\pi}{6},0\right)\) or awrt \((6.93,0)\), at the crossing of \(G_1\) and the initial line.} \\
\nopagebreak & & [3] & \\
7(c) & {\(r=6\sec\left(-\alpha+\tfrac{\pi}{6}\right)\)} & M1 & {Substitutes \(\theta=-\alpha\) into the equation of \(G_1\).} \\
& {\(\left(6\sec\left(\tfrac{\pi}{6}-\alpha\right),\,-\alpha\right)\)} & A1 & {Both coordinates, oe such as \(\frac{6}{\cos\left(\frac{\pi}{6}-\alpha\right)}\) or expanded.} \\
\nopagebreak & & [2] & \\
7(d) & {\(12\sqrt3=\tfrac12\left(\their{4\sqrt3}\right)\left(\their{6\sec\left(\tfrac{\pi}{6}-\alpha\right)}\right)\sin\alpha\)} & M1 A1 & {\msNote{M1}{Forms the triangle's area in terms of \(\alpha\), e.g. \(\frac12ab\sin C\) with their sides from \(O\) and angle \(\alpha\), or \(\frac12\int r^2\,\mathrm{d}\theta\), equal to \(12\sqrt3\).}\\ \msNote{A1}{Correct equation, oe.}} \\
& {\(\sin\alpha=\cos\left(\tfrac{\pi}{6}-\alpha\right)\)\\[4pt] \(\phantom{\sin\alpha}=\tfrac{\sqrt3}{2}\cos\alpha+\tfrac12\sin\alpha\)\\[4pt] \(\Rightarrow\tan\alpha=\ldots\)} & dM1 & {Expands the compound angle and rearranges to an equation in \(\tan\alpha\) alone. Condone sign slips.} \\
& {\(\tan\alpha=\sqrt3\Rightarrow\alpha=\tfrac{\pi}{3}\)} & A1 & {cso, with no other values given.} \\
\nopagebreak & & [4] & \\
\nopagebreak \SetCell[c=2]{r} Total for Question 7 & & 12 & \\
\end{msscheme}

\msQuestion{8}
\begin{msscheme}{
  hline{3-4,6-9,11-13,15-19} = {2-Z}{0.4pt, msGrid},
  hline{5,10,14,20} = {0.4pt, msGrid},
  row{20} = {font=\bfseries},
}
Question & Scheme & Marks & Notes \\
8(i) & {\(2(1+\cos 3\theta) = 0 \Rightarrow \cos 3\theta = -1\)\\[4pt] \(\phantom{2(1+\cos 3\theta)} \Rightarrow \theta = \ldots\)} & M1 & {Sets \(r = 0\) and solves \(\cos 3\theta = -1\) for at least one value of \(\theta\).} \\
& {\(\alpha = \frac{\pi}{3}\)} & A1 & {In terms of \(\pi\); 1.05 alone scores A0.} \\
\nopagebreak & & [2] & \\
8(ii)(a) & {\(f\left(\frac{2k\pi}{3}-\theta\right) = 2[1+\cos(2k\pi - 3\theta)]\)\\[4pt] \(\phantom{f} = 2[1+\cos 2k\pi\cos 3\theta + \sin 2k\pi\sin 3\theta]\)} & M1 & {Expands \(\cos(2k\pi - 3\theta)\) by the addition formula, or uses the evenness and period \(2\pi\) of cosine.} \\
& {\(\phantom{f} = 2(1+\cos 3\theta) = f(\theta)\)} & A1 & {Reaches \(f(\theta)\) without error for general integer \(k\), not a particular value.} \\
& {\(f\left(\theta+\frac{2\pi}{3}\right) = 2[1+\cos(3\theta + 2\pi)]\)\\[4pt] \(\phantom{f} = 2[1+\cos 3\theta\cos 2\pi - \sin 3\theta\sin 2\pi]\)} & M1 & {Expands \(\cos(3\theta + 2\pi)\) by the addition formula, or uses the period \(2\pi\) of cosine.} \\
& {\(\phantom{f} = 2(1+\cos 3\theta) = f(\theta)\)} & A1~AG & {Correct substituted expression leading without error to \(f(\theta)\), with \(= f(\theta)\) or a statement in words to conclude.} \\
\nopagebreak & & [4] & \\
8(ii)(b) & {Reflection in \(\theta = \beta\): \(\theta \to 2\beta - \theta\)\\[4pt] \(2\beta = \frac{2k\pi}{3} \Rightarrow \theta = \frac{k\pi}{3}\)} & B1 & {Lines \(\theta = \frac{k\pi}{3}\) oe, linked to the first identity.} \\
& {\(\theta = 0,\ \frac{\pi}{3},\ \frac{2\pi}{3}\)} & B1 & {All three and no others.} \\
& {Rotation through \(\frac{2\pi}{3}\)} & B1 & \\
\nopagebreak & & [3] & \\
8(iii) & \SetCell[c=3]{c} \msFigure{\begin{tikzpicture}[scale=1.1]
\def\radScale{2}
\def\frequency{3}
\pgfmathsetmacro{\maxRadius}{2*\radScale}
\pgfmathsetmacro{\axisLength}{1.2*\maxRadius}
\coordinate (O) at (0,0);
\coordinate (H) at (\axisLength,0);
\coordinate (P) at (0:\maxRadius);
\coordinate (Q) at (120:\maxRadius);
\coordinate (R) at (240:\maxRadius);
\foreach \centreAngle in {0,120,240} {
  \pgfmathsetmacro{\startAngle}{\centreAngle-180/\frequency}
  \pgfmathsetmacro{\endAngle}{\centreAngle+180/\frequency}
  \draw[thick, black]
    plot[smooth, samples=121, domain=\startAngle:\endAngle, variable=\t]
    ({\radScale*(1+cos(\frequency*\t))*cos(\t)},
     {\radScale*(1+cos(\frequency*\t))*sin(\t)});
}
\draw[thick, black, -{Stealth[length=3mm]}] (O) -- (H)
  node[right] {Initial line};
\foreach \point in {P,Q,R} {
  \fill[black] (\point) circle (1.5pt);
}
\node[above right] at (P) {$\left(4,0\right)$};
\node[above left] at (Q) {$\left(4,\frac{2\pi}{3}\right)$};
\node[below left] at (R) {$\left(4,\frac{4\pi}{3}\right)$};
\node[below left=2pt] at (O) {$O$};
\end{tikzpicture}} & & \\
& \msFeature{Shape} & B1 & {Three equal closed lobes meeting at \(O\), one centred on the initial line, the others in roughly the correct directions.} \\
& {\(r_{\max} = 4\)} & B1 & \\
& {\(\cos 3\theta = 1 \Rightarrow 3\theta = 2k\pi \Rightarrow \theta = \ldots\)} & M1 & {Sets \(\cos 3\theta = 1\) and finds at least one nonzero value of \(\theta\).} \\
& {\(\theta = 0,\ \frac{2\pi}{3},\ \frac{4\pi}{3}\)} & A1 & {All three, or 0, awrt 2.09 and 4.19, with no others in \(0 \leq \theta < 2\pi\).} \\
\nopagebreak & & [4] & \\
\nopagebreak \SetCell[c=2]{r} Total for Question 8 & & 13 & \\
\end{msscheme}

\msQuestion{9}
\begin{msscheme}{
  hline{3-4,6-11} = {2-Z}{0.4pt, msGrid},
  hline{5,12} = {0.4pt, msGrid},
  row{12} = {font=\bfseries},
}
Question & Scheme & Marks & Notes \\
9(a) & {\(r\cos\left(\theta-\frac{\pi}{4}\right)=3\)\\[4pt] \(\Rightarrow \frac{r\cos\theta+r\sin\theta}{\sqrt2}=3\)} & M1 & {Expands \(\cos\left(\theta-\frac{\pi}{4}\right)\) and replaces \(r\cos\theta\), \(r\sin\theta\) by \(x\), \(y\). Condone a slip in the value of \(\cos\frac{\pi}{4}\).} \\
& {\(x+y=3\sqrt2\)} & A1~AG & {Correct expansion with \(\cos\frac{\pi}{4}=\sin\frac{\pi}{4}=\frac{1}{\sqrt2}\), leading without error to the printed equation.} \\
\nopagebreak & & [2] & \\
9(b) & {\(\sin2\theta=2\cos^2\left(\theta-\frac{\pi}{4}\right)-1\)\\[4pt] \(\phantom{\sin2\theta}=\frac{18}{r^2}-1\)\\[4pt] \(r=5-2\left(\frac{18}{r^2}-1\right)\)} & M1 & {Eliminates \(\theta\) using \(L\) and a double-angle identity, e.g. via \((x+y)^2=r^2(1+\sin2\theta)\), to an equation in \(r\) alone. Condone sign slips.} \\
& {\(r^3-7r^2+36=0\)} & A1 & {oe, e.g. \(r=7-\frac{36}{r^2}\).} \\
& {\((r-3)(r-6)(r+2)=0 \Rightarrow r=\ldots\)} & dM1 & {Solves their cubic by any means, calculator included, to obtain values of \(r\).} \\
& {\(r=3,\ r=6\) (\(r=-2\) rejected as \(3\le r\le7\))} & A1 & {Both values, with \(r=-2\) rejected for any valid reason.} \\
& {\(\cos\left(\theta-\frac{\pi}{4}\right)=\frac{3}{\their{r}}\), \(-\frac{\pi}{2}<\theta-\frac{\pi}{4}<\frac{\pi}{2}\)\\[4pt] \(\Rightarrow \theta=\ldots\)} & M1 & {Substitutes a positive value of their \(r\) into either equation and solves for \(\theta\) in \(-\frac{\pi}{4}<\theta<\frac{3\pi}{4}\).} \\
& {\(\left(3,\frac{\pi}{4}\right),\ \left(6,-\frac{\pi}{12}\right),\ \left(6,\frac{7\pi}{12}\right)\)\\[4pt] No other solutions in the range of \(L\)} & A1 & {cso. All three exact points and no others, with a statement that the cosine equations have no further solutions in the range of \(L\).} \\
\nopagebreak & & [6] & \\
\nopagebreak \SetCell[c=2]{r} Total for Question 9 & & 8 & \\
\end{msscheme}

\msQuestion{10}
\begin{msscheme}{
  hline{3-5,7-9,11-15,17-20} = {2-Z}{0.4pt, msGrid},
  hline{6,10,16,21} = {0.4pt, msGrid},
  row{21} = {font=\bfseries},
}
Question & Scheme & Marks & Notes \\
10(a) & {\((r\cos\theta)^2-(r\sin\theta)^2=6r\cos\theta\)} & M1 & {Substitutes both into the Cartesian equation.} \\
& {\(r^2\cos2\theta=6r\cos\theta \Rightarrow r\cos2\theta=6\cos\theta\)} & A1 & \\
& {\(x>0 \Rightarrow \cos\theta>0\), \(r>0 \Rightarrow \cos2\theta>0\)\\[4pt] \(\Rightarrow -\frac{\pi}{4}<\theta<\frac{\pi}{4}\), \(r=6\cos\theta\sec2\theta\)} & A1~AG & {Printed equation reached without error, with the range deduced from \(\cos\theta>0\) and \(\cos2\theta>0\). A0 if the range is only stated.} \\
\nopagebreak & & [3] & \\
10(b)(i) & {\(0<\cos2\theta=\cos^2\theta-\sin^2\theta\leq\cos^2\theta\leq\cos\theta\)} & M1 & {Attempts a valid justification: bounds \(\cos2\theta\) by \(\cos\theta\) on the interval, sets \(\frac{\mathrm{d}r}{\mathrm{d}\theta}=0\), or writes \(r^2=2x^2-6x\) for \(x\geq6\).} \\
& {\(r=\dfrac{6\cos\theta}{\cos2\theta}\geq6\), equality only at \(\theta=0\)} & A1 & {Or an equivalent conclusion: the stationary point shown to be a minimum, or \(r^2\) increasing for \(x\geq6\).} \\
& {\(N=(6,0)\)} & A1 & \\
\nopagebreak & & [3] & \\
10(b)(ii) & {\(6\sqrt3\cos2\theta=6\cos\theta\)\\[4pt] \(\Rightarrow \sqrt3(2\cos^2\theta-1)=\cos\theta\)} & M1 & {Equates \(6\sqrt3\) to the polar equation and uses a double-angle formula to reach an equation in \(\cos\theta\) only; or eliminates \(y\) using \(x^2+y^2=108\).} \\
& {\(2\sqrt3\cos^2\theta-\cos\theta-\sqrt3=0\)} & A1 & {oe, or \(2x^2-6x-108=0\) oe.} \\
& {\((2\cos\theta-\sqrt3)(\sqrt3\cos\theta+1)=0 \Rightarrow \cos\theta=\ldots\)} & dM1 & {Solves their three-term quadratic by factorising (product of constants correct), the formula or completing the square; calculator allowed.} \\
& {\(\cos\theta=\frac{\sqrt3}{2}\) (\(-\frac{1}{\sqrt3}\) rejected as \(\cos\theta>0\))} & A1 & {Other root rejected with a reason, eg \(\cos\theta>0\) on the interval; or \(x=9\) with \(x=-6\) rejected.} \\
& {\(P\left(6\sqrt3,\frac{\pi}{6}\right)\), \(Q\left(6\sqrt3,-\frac{\pi}{6}\right)\)} & A1 & {Both points exact, angles in the stated interval; \(\frac{11\pi}{6}\) for \(Q\) scores A0.} \\
\nopagebreak & & [5] & \\
10(b)(iii) & {\(NP^2=6^2+(6\sqrt3)^2-2(6)(6\sqrt3)\cos\frac{\pi}{6}\)} & M1 & {Uses the cosine rule in triangle \(ONP\) with \(ON=6\), their \(OP\) from (b)(ii) and angle \(\frac{\pi}{6}\), or the distance formula, or finds the Cartesian coordinates of \(P\).} \\
& {\(NP=6\)} & A1 & {\(NP=6\), or \(P(9, 3\sqrt3)\) oe.} \\
& {\(\angle ONP=\pi-2\left(\frac{\pi}{6}\right)\ [=\tfrac{2\pi}{3}]\)\\[4pt] \(\angle PNQ=2\pi-2\angle ONP=\ldots\)} & M1 & {Uses \(ON=NP\) to find \(\angle ONP\) and symmetry about the initial line to reach \(\angle PNQ\); or any valid route, eg the direction of \(NP\).} \\
& {\(\angle PNQ=\frac{2\pi}{3}\)} & A1 & {Exact, in radians; \(120^\circ\) scores A0.} \\
\nopagebreak & & [4] & \\
\nopagebreak \SetCell[c=2]{r} Total for Question 10 & & 15 & \\
\end{msscheme}

\end{document}
